\section{Prodotti e somme in \(\ctSet\)}
Le operazioni insiemistiche di somma e prodotto sono già state impiegate diverse volte nel capitolo \ref{chap_cat_fun_nat} in una versione intuitiva, che dopotutto è quella con cui sono state introdotte in algebra o analisi; lo scopo di questa prima sezione è `ristrutturarle' in prospettiva categoriale, e descriverle come primi importanti esempi di costruzioni universali. Nel dover dimostrare le proprietà di unicità per insiemi e funzioni, ci avvarremo del principio di estensionalità per le funzioni, che è un assioma di ZF(C) (dunque implicito nelle nostre assunzioni di base): due funzioni	\(f,g \colon A\to B\) sono uguali se e solo se sono uguali le immagini \(f(a),g(a)\) per ogni \(a\in A\). In questo modo, quando si dovrà dimostrare che una certa funzione di insiemi \(f \colon A\to B\) è l'unica con una certa proprietà \(\mathcal{P}\), dimostreremo che, se \(g \colon A \to B\) è un'altra funzione con la proprietà \(\mathcal{P}\), allora \(f(a)=g(a)\) per ogni \(a\in A\).
\subsection{Prodotti in \(\ctSet\)}\label{prod_in_Set}
Per produrre un \emph{prodotto Cartesiano} \(A\times B\), viene costruito un insieme di coppie ordinate di elementi presi da due \emph{fattori}  \(A,B\). Come ogni insieme, \(A\times B\) è caratterizzato da quali sono i suoi elementi. Le regole di `introduzione' di un tipo/insieme permettono di costruire un suo elemento a partire da dati più semplici; le regole di `eliminazione' lo destrutturano nelle sue componenti; una o più regole di `computazione' prescrivono il modo in cui introduzione ed eliminazione si annullano a vicenda.
\begin{hExample}[Regole di introduzione ed eliminazione per il prodotto Cartesiano]{fund}\label{intro_elim_prodotto}
	Dati due insiemi \(A\) e \(B\), il loro \emph{prodotto Cartesiano} \(A\times B\) consiste nell'insieme di tutte le coppie ordinate di elementi di \(A\) e di \(B\):\footnote{Il prodotto è chiamato così in onore del filosofo francese Pierre Descartes (\(*\)1596-\(\dag\)1650) che per primo si accorse (si dice, osservando il soffitto dove si era posata una mosca) che un punto del piano può essere univocamente determinato da una coppia di coordinate \((x,y)\in\bbR\times\bbR\).}
	\[
		A\times B= \{\pair ab \mid a\in A, b\in B\}
	\]
	Per chi è familiare con questa nomenclatura, \(A\times B\) è il `tipo' i cui `termini' sono destrutturabili come coppie: le seguenti `frazioni'
	\[\frac{t \in A\times B}{\pi_A\, t \in A}\qquad \frac{t \in A\times B}{\pi_B\, t \in B}\qquad \frac{a \in A\quad b \in B}{\pair ab \in A\times B}\qquad \frac{t \in A\times B}{t = \pair{\pi_A\, t}{\pi_B\, t}}\]
	si leggono come deduzioni (supponendo che <numeratore>, allora <denominatore>), e dicono quattro cose:
	\begin{itemize}
		\item ogni volta che viene dato un elemento \(t \in A\times B\), si ottiene un elemento \(\pi_A t\) la `prima componente di \(t\)';
		\item ogni volta che viene dato un elemento \(t \in A\times B\), si ottiene un elemento \(\pi_B t\) la `seconda componente di \(t\)';
		\item ogni volta che vengono dati due elementi \(a \in A, b \in B\), si ottiene un elemento \(\pair ab \in A\times B\);
		\item ogni elemento \(t \in A\times B\) è della forma \(\pair{\pi_A\, t}{\pi_B\, t}\).
	\end{itemize}
	L'ultima di queste condizioni a sua volta dice due cose:
	\begin{itemize}
		\item in \(A\times B\) non ci sono che elementi della forma \(\pair ab\) per qualche \(a \in A\) e \(b \in B\);
		\item le uguaglianze in \ref{}, \ref{}, \ref{} valgono per `lambda-astrazione' di \(t\), cioè
		      \[\pair{\pi_A}{\pi_B} = \id_{A\times B} \]
	\end{itemize}
\end{hExample}
In questi termini, il prodotto è descritto come una operazione di costruzione di un insieme a partire da due dati: questo è ragionevole, dato che abbiamo già visto in \ref{} che esiste un funtore
\[\dmFun{\blank\times\blank}{\ctSet\times\ctSet}{\ctSet.}\]
Come già detto, il piano Cartesiano ortogonale della geometria analitica può essere identificato con \(\bbR^2=\bbR\times\bbR\), l'insieme delle coppie \emph{ordinate} (questa condizione è cruciale) di numeri reali.

L'esempio \ref{} lascia poi intendere che ci si debba più correttamente riferire al prodotto \(A\times B\) come la terna
\[
	A\times B\,,\quad \pi_A\colon A\times B\to A\,,\quad \pi_B\colon A\times B\to B\,,
\]
specificata insieme alle funzioni \(\pi_A,\pi_B\) di \emph{proiezione} sui due fattori \(A\) e \(B\), definite in \ref{}.

Anche il prodotto Cartesiano di insiemi può essere caratterizzato da una proprietà che non riguarda direttamente i suoi elementi, ma li caratterizza in termini delle funzioni verso \(A\times B\) (gli elementi di \(A\times B\) essendo i \emph{punti} della forma \(\pair ab \colon \singleton \to A\times B\)). Per chiarire questo aspetto, consideriamo una funzione che ha \(A\times B\) come codominio, cioè della forma
\[\label{prodotti_h}
	h\colon X\to A\times B
\]
per qualche insieme \(X\). Essa può essere composta con le proiezioni \(\pi_A,\pi_B\) per ottenere due funzioni
\[\label{prodotti_f_e_g}
	\pi_A\cmp h\colon X\to A \qquad \pi_B\cmp h\colon X\to B
\]
che determinano univocamente \(h\). Infatti, si osservi come
\begin{itemize}
	\item grazie a \ref{eq_comprule} l'immagine di \(h(x)\) può essere \emph{ridotta} alla coppia ordinata
	      \[\pair{(\pi_A\cmp h)(x)}{(\pi_B\cmp h)(x)};\]
	\item comunque siano date due funzioni \(f \colon X\to A\) e \(g \colon X \to B\), è possibile \emph{costruire} la coppia ordinata \(\pair{f(x)}{g(x)}\) e la funzione \(f\pmap g\) che abbiamo già introdotto in \ref{puniv_prodotto_naturale} che ne risulta è l'immagine di \(\pair fg\), nell'insieme a sinistra, rispetto alla biiezione naturale (già dimostrata in \ref{})
	      \[\Hom{\ctSet}(X,A\times B)\cong\Hom{\ctSet}(X,A)\times\Hom{\ctSet}(X,B).\]
\end{itemize}
Queste due corrispondenze di riduzione e costruzione, o introduzione ed eliminazione, che sono inverse una dell'altra, stanno alla base del ragionare mediante proprietà universali.
\begin{remark}[Proprietà universale del prodotto Cartesiano]\label{rmk_prodotto}
	Questa proprietà è `caratterizzante' nel senso che segue: una terna costituita da un insieme \(P\) e due funzioni \(p_1,p_2\) del tipo
	\[
		(P,\,  p_1\colon P\to A,\,  p_2\colon P\to B)
	\]
	si dice un \emph{prodotto Cartesiano} di \(A\) e \(B\), se, e solo se, per ogni insieme \(X\) e per ogni coppia di funzioni \(f \colon X \to A\) e \(g \colon X \to B\) come in \eqref{prodotti_f_e_g}, esiste un'unica funzione \(f\pmap g \colon X \to P\) come in \eqref{prodotti_h} tale che
	\[
		p_1\cmp (f\pmap g) = f\,,\qquad p_2\cmp (f\pmap g)=g.
	\]
\end{remark}
In effetti, è immediato verificare che \(A\times B\) e le proiezioni canoniche \(\pi_A,\pi_B\) soddisfano questa proprietà. Ma se \(P\) soddisfa questa proprietà, possiamo dimostrare quanto segue.
\begin{lemma}\label{prodotto_unico_upto_iso}
	Siano \((P,p_1,p_2)\) e \((Q,q_1,q_2)\) due prodotti	cartesiani di \(A\) e \(B\). Allora esiste un unico isomorfismo \(u= p_1\pmap p_2 \colon P\to Q\) di inversa \(u^{-1} = q_1\pmap q_2\). In particolare, se \((P,p_1,p_2)\) è un prodotto Cartesiano di \(A\) e \(B\) come in \ref{rmk_prodotto}, allora esiste un'unica biiezione \(P\cong A\times B\), che rende \(P\) `essenzialmente' composto dalle coppie \(\pair ab\) come in \ref{}.
\end{lemma}
Si osservi di nuovo il ruolo dell'aggettivo \emph{unica}: fissate le terne \((P,p_1,p_2)\) e \((Q,q_1,q_2)\) (e solo allora), la biiezione possibile tra \(P\) e \(Q\) è \emph{solo una}. In termini più semplici, ci possono essere biiezioni tra \(P\) e \(Q\) che non rendono commutativi i diagrammi di \ref{rmk_prodotto}, ma ce n'è una sola che li rende commutativi, e questa è quella cui ci si riferisce nelle discussioni su oggetti universali. Identiche considerazioni valgono per i prodotti	in categorie astratte, e per tutti gli altri oggetti universali che incontreremo.
\begin{proof}
	Dato che \(Q\) è un prodotto, esiste un'unica	funzione \(u \colon P \to Q\) che fa commutare i due triangoli del diagramma
	\[\xymatrix{
			& P \ar@{.>}[d]^{u} \ar[dl]_{p_1} \ar[dr]^{p_2}\\
			A & Q \ar[r]_{q_2}\ar[l]^{q_1}& B
		}\]
	Di converso, dato che	\(P\) è un prodotto, esiste un'unica funzione \(v \colon Q \to P\) che fa commutare i due triangoli del diagramma
	\[\xymatrix{
			& Q \ar@{.>}[d]^v \ar[dl]_{q_1} \ar[dr]^{q_2}\\
			A & P \ar[r]_{p_2}\ar[l]^{p_1}& B.
		}\]
	Del resto, \(v\) deve essere l'inversa di \(u\), perché le composizioni \(v\cmp u\) e \(u\cmp v\) fanno commutare i due triangoli dei diagrammi
	\[\xymatrix{
			& P \ar@{.>}[d]^{vu} \ar[dl]_{p_1} \ar[dr]^{p_2}\\
			A & P \ar[r]_{p_2}\ar[l]^{p_1}& B
		}\qquad\qquad \xymatrix{
			& Q \ar@{.>}[d]^{uv} \ar[dl]_{q_1} \ar[dr]^{q_2}\\
			A & Q \ar[r]_{q_2}\ar[l]^{q_1}& B.
		}\]
	proprietà che esse hanno in comune con	le identità di \(P\) e \(Q\); dunque, per unicità, si ha \(v\cmp u = \id_P\) e \(u\cmp v = \id_Q\).
\end{proof}
Si osservi che in particolare, istanziando la proprietà di \ref{rmk_prodotto} quando \(X = \singleton\) è un insieme con un solo elemento, le funzioni \(g\colon \singleton\to B\), \(f\colon \singleton\to A\) scelgono rispettivamente un elemento \(g(\bullet) = b \in B\) e \(f(\bullet) = a \in A\); la funzione \(f\pmap g\colon \singleton\to P\) sceglie a sua volta un elemento \((f\pmap g)(\bullet) = t \in P\), che però ora deve essere della forma \(\pair{p_1 t}{p_2 t}\).

La proprietà di \ref{} si traduce allora in: per ogni coppia di elementi \(a\in A\) e \(b\in B\), esiste (suriettività) un unico (iniettività) elemento \(c\in P\) tale che \(p_1 c=a\) e \(p_2 c=b\). Questo è un modo alternativo di stabilire una biiezione \(\tau \colon P\to A\times B\) soggetta alla proprietà addizionale di commutare con le proiezioni, cioè tale che \(\pi_A\cmp \tau =p_1\) e \(\pi_B\cmp \tau =p_2\). La biiezione è poi unica, ossia non solo è possibile identificare un prodotto Cartesiano \(P\) ad un insieme di coppie, ma questa identificazione si fa in un unico modo.

In sintesi, abbiamo enunciato e dimostrato la
\begin{definition}\label{pruniv_prodotto_tfae}
	Siano \(A,B\) due insiemi. Il \emph{prodotto Cartesiano} \(A\times B\) di \(A\) e \(B\) è un insieme \(A\times B\) con due funzioni di \emph{proiezione}
	\[\xymatrix{
			A & A\times B \ar[r]\ar[l]& B
		}\]
	tali che una di queste proprietà, equivalenti tra loro, è soddisfatta:
	\begin{itemize}
		\item (proprietà universale diagrammatica) per ogni spanna
		      \[\xymatrix{
				      A & A\times B \ar[r]\ar[l]& B
			      }\]
		      esiste un unico omomorfismo di spanne denotato \(\pair fg\)	che rende commutativo il diagramma
		      \[\xymatrix{
			      & X \ar@{.>}[d]|{\pair fg} \ar[dl]_f \ar[dr]^g\\
			      A & A\times B \ar[r]_{p_2}\ar[l]^{p_1}& B
			      }\]
		\item (proprietà universale mediante rappresentabilità) Per ogni \(X\), la funzione
		      \[\xymatrix@R=0cm{
			      \varpi_X : \Hom{\ctSet}(X,A\times B) \ar[r]& \Hom{\ctSet}(X,A)\times\Hom{\ctSet}(X,B)\\
			      h \ar@{|->}[r]& (\pi_A\cmp h, \pi_B\cmp h)
			      }\]
		      è una biiezione naturale in \(X\) (cioè, l'insieme \(A\times B\) rappresenta il funtore \(\Hom{\ctSet}(\blank,A)\times\Hom{\ctSet}(\blank,A) \colon \ctSet^\op \to \ctSet\) che manda \(X\) in \(\Hom{\ctSet}(X,A)\times\Hom{\ctSet}(X,B)\)).
	\end{itemize}
\end{definition}
\begin{remark}\label{tediosi_ma_anche_no}
	È probabile che questi argomenti risultino persino troppo pedanti per chi legge; tuttavia, ragionare in questi termini è il cuore della teoria delle proprietà universali, che `si dimostrano tutte allo stesso modo'. Ancorare la tecnica di dimostrazione a un esempio estremamente banale è uno strumento pedagogico che permette di non oscurare la tecnica	con i dettagli specifici di una costruzione più complessa (che andrebbe capita e motivata\dots).
\end{remark}
Altrettanto probabilmente, chi legge avrà intuito che c'è una somiglianza tra il modo di ragionare appena descritto, e il TFPU di \ref{tfpu}, e ancor prima di \ref{init_ha_solo_id}; questa non è una coincidenza. La dimostrazione del seguente fatto consiste nel notare che le due condizioni di \ref{rmk_prodotto} sono esattamente le condizioni che rendono \(P\) un oggetto terminale.
\begin{proposition}\label{prodotto_terminal_span}
	Si rammenti la definizione della categoria \(\ctSpan(A,B)\) in \ref{}, delle spanne di insiemi, che ha
	\begin{itemize}
		\item per oggetti i diagrammi \(\xymatrix{ A & X \ar[r]^-r \ar[l]_-l & B }\);
		\item per frecce le funzioni \(f \colon X\to Y\) che rendono commutativi i diagrammi
		      \[\xymatrix{
			      X\ar[dr]_f \ar[r]^-l\ar[d]_r & A \\
			      B & Y. \ar[l]^-{r'}\ar[u]_-{l'}
			      }\]
	\end{itemize}
	Un oggetto \(\xymatrix{ A & P \ar[r]^-{p_B}\ar[l]_-{p_A} & B }\) è allora terminale in \(\ctSpan(A,B)\) se e solo se \((P,p_A,p_B)\) è un prodotto Cartesiano di \(A\) e \(B\).
\end{proposition}
Siamo pronti per enunciare la proprietà universale del prodotto (di due oggetti) in una categoria \(\ctC\); seguiranno molti esempi di prodotti binari nelle categorie definite nel capitolo precedente, così come diversi controesempi che mostrano come alcune categorie non abbiano prodotti, o non li abbiano per ogni coppia di oggetti. La sezione \ref{} si occuperà di tracciare le prime conseguenze della definizione, e le sue proprietà elementari.
\begin{definition}[Prodotto binario in \(\ctC\)]\label{def_prodotto_in_una_cat}
	Sia \(\ctC\) una categoria e \(A,B\in\ctC_0\) due oggetti. Un oggetto \(P\) di \(\ctC\) è un prodotto di \(A\) e \(B\) se una delle seguenti condizioni equivalenti è soddisfatta:
	\begin{enumtag}{pd}
		\item\label{pd_1} per ogni diagramma \(A \xot{z_A} Z \xto{z_B} B\), esiste un'unica freccia \(\pair{z_A}{z_B} \colon Z\to P\)	che rende commutativi i due triangoli del diagramma
		\[\xymatrix{
			& Z \ar@{.>}[d]|{\pair{z_A}{z_B}} \ar[dl]_{z_A}\ar[dr]^{z_B}\\
			A & P \ar[r]^{p_2}\ar[l]_{p_1} & B
			}\]
		\item\label{pd_2} la spanna \(\xymatrix{
		A & P \ar[r]^-{p_B}\ar[l]_-{p_A}& B
		}\) è un oggetto terminale in \(\ctSpan_\ctC(A,B)\) (la categoria delle spanne di \(\ctC\) tra \(A\) e \(B\) definita in \ref{});
		\item\label{pd_3} l'oggetto	\(P\) rappresenta (nel senso di \ref{}) il funtore \(\varpi_{AB}=\Hom{\ctC}(\blank,A)\times\Hom{\ctC}(\blank,B) \colon \ctC^\op \to \ctSet\) che manda \(X\) in \(\Hom{\ctC}(X,A)\times\Hom{\ctC}(X,B)\).
	\end{enumtag}
\end{definition}
\begin{remark}[Alcune osservazioni sulla nozione generale di prodotto]\label{rmk_prodotto_gen}
	Se una qualsiasi di queste condizioni	è soddisfatta, si dice che \(P\) è un prodotto di \(A\) e \(B\), e si denota con \(A\times B\). Dire `\(P\) è un prodotto' in effetti è una denominazione ambigua; il prodotto è, per definizione, \emph{l'intera spanna} \((P,p_A,p_B)\), e l'oggetto \(P\) ne è l'apice. Ma è consuetudine ammettere questo lieve abuso di terminologia. In questo modo, quando si dice che `\(P\) è un prodotto', solitamente si intende che le braccia della spanna \(p_A \colon P\to A\) e \(p_B \colon P \to B\) esistono ed è ovvio determinarle dal contesto del problema.

	Allora diremo che \(\ctC\) `ha prodotti' o `è una categoria con prodotti' se per ogni \(A,B\in\ctC_0\) esiste il prodotto \(A\times B\).

	Si osservi che la condizione \ref{pd_3} sembra non fare menzione della spanna di apice \(P\); in realtà però la sua esistenza è implicata dalla condizione di rappresentabilità: se \(P\) rappresenta il funtore \(\Hom{\ctC}(\blank,A)\times\Hom{\ctC}(\blank,B)\), allora esiste una biiezione naturale
	\[\alpha_P : \Hom{\ctC}(P,P)\cong\Hom{\ctC}(P,A)\times\Hom{\ctC}(P,B)\]
	e l'immagine di \(\id_P\) mediante \(\alpha_P\) consta di due frecce \(p_1 \colon P\to A\) e \(p_2 \colon P\to B\), che rendono \(P\) un prodotto di \(A\) e \(B\) secondo la prima condizione. Questa è semplicemente una maniera di rifrasare il fatto che l'oggetto \((P,p_1,p_2)\) è terminale nella categoria degli elementi di \(\varpi_{AB}\), come definita in \ref{}, che coincide esattamente con \(\ctSpan_\ctC(A,B)\).

	Considerazioni del tutto analoghe saranno valide per la nozione generale di limite, che sarà definita in \ref{}; considerazioni del tutto analoghe, dualizzate, saranno valide per la nozione generale di colimite, che sarà definita in \ref{}.
\end{remark}
\subsection{Prodotti in altre categorie `concrete'}
I primi corsi di algebra hanno presentato diverse classi di strutture algebriche come monoidi, gruppi, anelli, e a volte reticoli algebrici, grafi, etc.; in ciascuna di queste strutture esiste una nozione di prodotto tra due oggetti, che è caratterizzata da una proprietà universale del tutto analoga a quella del prodotto Cartesiano di insiemi. Cioè: il prodotto di due gruppi (o anelli, spazi vettoriali, reticoli, grafi\dots) \(G\) e \(H\) si definisce ponendo sul prodotto Cartesiano \(G\times H\) degli insiemi sottostanti una operazione componente per componente, non perché questa definizione sia ovvia, ma perché è l'unica che soddisfa la giusta proprietà universale.
\begin{hExample}[Prodotti in categorie di insiemi strutturati]{fund}\label{prodotti_cat_strutturate}
	Raccogliamo vari esempi, tutti formalmente analoghi, di prodotti in categorie di insiemi con struttura.
	\begin{itemize}
		\item Il prodotto di due gruppi \(G,H\) si costruisce ponendo sul prodotto Cartesiano \(G\times H\) la moltiplicazione definita componente per componente,
		      \[(g,h) \cdot_{G\times H} (g',h') \defeq (g\cdot_G g', h\cdot_H h');\]
		      l'identità è ovviamente \((1_G,1_H)\), l'inverso di \((g,h)\) è \((g^{-1},h^{-1})\). Similmente si procede per definire il prodotto di due gruppi abeliani, di due spazi vettoriali \(V,W\), di due anelli, et cetera, fino al punto che è facile sospettare ci sia un teorema generale che afferma che in ogni categoria di modelli di una segnatura algebrica (come in \ref{}), il prodotto di due modelli	\(A\) e \(B\) si costruisce ponendo sul prodotto Cartesiano (fatto in \(\ctSet\)) dei loro supporti, una struttura definita componente per componente. Questo è vero.
		\item \Todo{}
		\item \Todo{}
	\end{itemize}
\end{hExample}
\begin{hExample}[Prodotti nella categoria degli spazi topologici]{fund}\label{ex_prodotti_in_top}
\end{hExample}
\begin{hExample}[Prodotti in categorie di insiemi ordinati]{fund}\label{ex_prodotti_in_pos}
	L'esempio degli insiemi ordinati è particolarmente interessante, perché mostra che ci sono due maniere di porre sul prodotto Cartesiano \(A\times B\) di due ordini parziali \((A,\le), (B,\preceq)\) un ordine;
	\begin{itemize}
		\item uno è l'ordine prodotto: \((a,b)\le_\times (a',b')\) se e solo se \(a\le a'\) in \(A\), e \(b\preceq b'\) in \(B\);
		\item uno è l'ordine lessicografico: \((a,b)\le_\text{lx} (a',b')\) se e solo se vale
		      \[(a\le a') \text{ oppure } \bigl( (a=a') \land (b\preceq b') \bigr).\]
	\end{itemize}
	La prima definizione è quella `giusta' dato che soddisfa la proprietà universale come enunciata in \ref{}; la seconda no --e quindi, in particolare, sebbene esista una ovvia biiezione monotòna \((A\times B,\le_\times) \cong (A\times B,\le_\text{lx})\), i due ordini non sono isomorfi.

	Ad esempio, se \(A=B=\{0\le 1\}\) si ha che \((A\times B,\le_\times)\) è l'ordine a quattro elementi isomorfo a \(P[2]\) (come in \ref{fig_n_cubo}), mentre \((A\times B,\le_\text{lx})\) è l'ordine totale \(\{(0,0) \le (0,1) \le (1,0) \le (1,1)\}\). Come è semplice verificare nel caso in cui \(B=\{x<y<z\}\) poi, la proiezione \(\pi_B \colon A\times B \to B\) sul secondo fattore non è monotòna.
\end{hExample}
\begin{hExample}[Prodotti in categorie della forma \(\presh\ctC\)]{fund}\label{ex_prodotti_in_presh}
	\Todo{e in particolare nelle categorie di funtori che hanno una presentazione `concreta', come i grafi diretti.}
\end{hExample}
La proprietà universale del prodotto è una maniera molto utile di determinare `isomorfismi canonici' tra oggetti: un esempio semplice ma illustrativo è il seguente.
\begin{lemma}\label{inrel_cop_isprod}
	Dati due insiemi \(A,B\), esiste una spanna di insiemi
	\begin{equation}
		\label{pow_cop_isprod}\xymatrix{
			\pow A & \ar[r]^{p_B}\ar[l]_{p_A}\pow{A+B} & \pow B
		}
	\end{equation}
	determinata dalle funzioni \(p_A= \pow{\iota_A}\) e \(p_B=\pow{\iota_B}\), dove \(\iota_A \colon A \to A+B\) e \(\iota_B \colon B \to A+B\) sono le inclusioni nell'unione disgiunta di insiemi. Esplicitamente, se \(U\subseteq A+B\),  \(p_A(U)= A \cap U\) e \(p_B(U)= B \cap U\).

	La spanna così definita soddisfa la proprietà universale del prodotto, e dunque esiste un `isomorfismo canonico'\footnote{Chi legge potrà trovare con piacere che questo implica che le cardinalità di questi insiemi sono uguali, e dunque che le `proprietà delle potenze' imparate alle scuole medie valgono in forza del fatto che il diagramma in \eqref{pow_cop_isprod} soddisfa la proprietà universale del prodotto. In particolare, non c'è niente di speciale nell'insieme `2' e analoghi isomorfismi del tipo \(X^{A+B}\cong X^A\times X^B\) valgono per ogni insieme \(X\).} tra \(\pow{A+B}\) e \(\pow A \times \pow B\), cioè una biiezione (che, in più, si può mostrare essere naturale in \(A,B\); si veda \ref{iso_aritmetici})
	\[\pow{A+B} \cong \pow A \times \pow B.\]
\end{lemma}
\begin{proof}
	La dimostrazione consiste nella verifica della proprietà universale per la spanna in \eqref{pow_cop_isprod}: dati un insieme \(X\) e una qualsiasi altra spanna
	\[\xymatrix{
			\pow A & \ar[r]^{f_B}\ar[l]_{f_A} X & \pow B
		}\]
	è possibile definire una funzione \(X \to \pow{A+B}\)	che rende commutativo il diagramma
	\[\xymatrix{
		&X\ar@{.>}[d]\ar[dr]^{f_B}\ar[dl]_{f_A}&\\
		\pow A & \ar[r]_-{p_B}\ar[l]^-{p_A} \pow{A+B} & \pow B
		}\]
	La funzione \(\pair{f_A}{f_B}\) è definita da \(\iota_A(f_A(x)) + \iota_B(f_B(x))\) per ogni \(x\in X\). La verifica che questa funzione rende commutativo il diagramma è semplice, e la lasciamo a chi legge. Per mostrare he \(\pair{f_A}{f_B}\) è l'unica definizione possibile, si renda precisa questa idea: se \(r \colon X \to \pow{A+B}\), allora \(r(x) = E+F\) per un'unica coppia di insiemi \(E\subseteq A,F\subseteq B\) (usando \ref{}); allora se \(p_A(r(x))=	f_A(x)\) e \(p_B(r(x)) = f_B(x)\), non può che essere \(p_A(r(x)) = A\cap E=E=f_A(x)\) e \(p_B(r(x)) = B\cap F=F=f_B(x)\), da cui \(r(x) = \iota_A(f_A(x)) + \iota_B(f_B(x))\).

	Concludiamo osservando due cose:
	\begin{itemize}
		\item Nella terminologia di \ref{}, diremo che il funtore \(\pow{(\blank)}\) `preserva' oppure che `commuta con' i prodotti (ma dato che è controvariante, manda \emph{coprodotti}, definiti in \ref{}, in prodotti).
		\item Il diagramma in \eqref{pow_cop_isprod} si può estendere a
		      \[\xymatrix{
			      \pow A\ar@<-.3em>[r]_{(\iota_{B})_*} & \ar@<.3em>[r]^-{p_B}\ar@<-.3em>[l]_-{p_A}\pow{A+B} & \pow B\ar@<.3em>[l]^-{(\iota_{A})_*}
			      }\]
		      dove abbiamo definito \((\iota_{A})_* (U) = \iota_A(U)\) e \((\iota_{B})_* (V) = \iota_B(V)\); quindi tra queste frecce valgono le relazioni
		      \[p_A((\iota_{A})_* (U)) = A\cap\iota_A(U) = \iota_A(A\cap U)=U,\quad
			      p_B((\iota_{B})_* (V)) = B\cap\iota_B(V) = \iota_B(B\cap V)=V\]
		      cosicché \(p_A,p_B\) sono epi spezzanti (risp., \((\iota_A)_*, (\iota_B)_*\) sono mono spezzanti). \qedhere
	\end{itemize}
\end{proof}
\begin{example}[Prodotti in \(\ctRel\)]\label{ex_prodotti_in_rel}
	Nella categoria \(\ctRel\) di insiemi e relazioni, come definita  in \ref{ex_cat_rels}, ogni coppia di oggetti \(A,B\) ammette un prodotto, \emph{ma}	questo prodotto non è il prodotto Cartesiano di insiemi, bensì l'insieme \(A + B\) (unione disgiunta di \(A\) e \(B\)), con la spanna di relazioni
	\[\xymatrix{
		A & A+B \ar[r]|@-{|}^-{p_B}\ar[l]|@-{|}_-{p_A}& B
		}\]
	definita (usando \ref{diff_pres_cat_rel}, la spanna è formata da funzioni \(\pow A \xot{p_A} A+B \xto{p_B} \pow B\)) come segue:
	\[p_A(t)=\begin{cases}
			\{t\}     & t\in A \\
			\emptyset & t\in B
		\end{cases}\quad
		p_B(t)=\begin{cases}
			\emptyset & t\in A \\
			\{t\}     & t\in B
		\end{cases}\]
\end{example}
\begin{example}[Prodotti di pseudoinsiemi]\label{ex_prodotti_in_bishop}
	Si ricordi da \ref{} che uno pseudoinsieme, o insieme di Bishop, è una coppia \((X,\sim)\) dove \(X\) è un insieme	e \(\sim\) è una relazione di equivalenza su \(X\). Il prodotto tra insiemi di Bishop	\((X,\sim_X)\) e \((Y,\sim_Y)\) è l'insieme di coppie \(X\times Y\) con la relazione di equivalenza definita dal prodotto	di relazioni di equivalenza, cioè
	\[(x,y)\sim_{X\times Y} (x',y') \defeq \bigl((x\sim_X x') \land (y\sim_Y y')\bigr).\]
	La proprietà universale è soddisfatta: se \(g \colon (Z,\approx)\to (Y,\sim_Y)\) e \(f \colon (Z,\approx)\to (X,\sim_X)\) sono funzioni che preservano le rispettive relazioni di equivalenza, allora la funzione \(\pair fg \colon Z\to X\times Y\) definita da \(\pair fg(z) = (f(z),g(z))\) è l'unica tale che
	\[z\approx z' \Rightarrow \bigl(f(z)\sim_X f(z') \land g(z)\sim_Y g(z')\bigr).\]
\end{example}
\begin{example}[Prodotti nell'incubo \(\ctSet/A\), e nella succuba \(A/\ctSet\)]\label{ex_prodotti_in_succuba_incubo}
	Nella categoria succuba sopra \(A\in\ctSet_0\), due oggetti \(f	\colon A\to X\) e \(g \colon A\to Y\) ammettono un prodotto, determinato dal prodotto Cartesiano \(X\times Y\) dei loro codomini, e dalla funzione \(\pair fg \colon A\to X\times Y\colon a \mapsto (f(a),g(a))\) per ogni \(a\in A\). Dunque, in \(A/\ctSet\), il prodotto di \(f \colon A\to X\) e \(g \colon A\to Y\) viene `creato' a partire dal prodotto in \(\ctSet\) dei codomini di \(f\) e \(g\).

	Per contro, nella categoria incubo sopra \(A\), due oggetti \(f \colon X\to A\) e \(g \colon Y\to A\) la situazione è irriducibile a una costruzione già nota. Il prodotto di due funzioni, viste come oggetti di \(\ctSet/A\), è il \emph{prodotto fibrato} sopra \(A\), che si costruisce come sottoinsieme (quasi sempre proprio) di \(X\times Y\) formato dalle coppie \((x,y)\) tali che \(f(x)=g(y)\), con le restrizioni delle proiezioni canoniche \(p_X|_{\pull XYA} \colon \pull XYA \subseteq X\times A\to X\colon (x,y)\mapsto x\) e \(p_Y|_{\pull XYA} \colon \pull XYA \subseteq X\times A\to Y\colon (x,y)\mapsto y\). Solo dopo questa restrizione la proprietà universale è soddisfatta, perché una spanna
	\[\xymatrix{
			(X,f) & (Z,h) \ar[r]^-r\ar[l]_-l & (Y,g)
		}\]
	si scompone a un diagramma
	\[
		\xymatrix{
		Z \ar@/^1pc/[rrd]^{r} \ar@/_1pc/[rdd]_{l} & & \\
		& \pull XYA \ar[r]^-{p_Y} \ar[d]_{p_X} & Y \ar[d]^{g} \\
		& X \ar[r]_{f} & A
		}
	\] con la proprietà che \(g\cmp r = f\cmp l=h\); ma allora la funzione \(\pair lh \colon Z\to X\times Y\) fattorizza sul sottoggetto \(\pull XYA\) come è stato definito, perché questa condizione è precisamente necessaria e sufficiente affinché \((l(z),r(z)) \in \pull XYA\).

	Sarà la sezione sui prodotti fibrati a chiarire la questione in modo definitivo.
\end{example}
\begin{example}[Prodotti in una categoria-poset]\label{ex_prodotti_in_catcomepos}
	\Todo{}
	Ad esempio,
	\begin{itemize}
		\item nella categoria-poset \(\pow{S}\), insieme delle parti di un insieme \(S\), il prodotto di \(A,B\subseteq S\) è l'intersezione \(A\cap B\) con le inclusioni \(A\supseteq A\cap B\subseteq B\);
		\item nella categoria-poset \((\bfN,\le)\), il prodotto di \(n,m\in\bfN\) è il loro minimo \(\min\{n,m\}\); mentre nella categoria-poset \((\bfN,\div)\), il prodotto di \(n,m\in\bfN\) è il loro massimo comun divisore \(\gcd(n,m)\).
	\end{itemize}
\end{example}
\begin{example}[Prodotti in una categoria-monoide]\label{ex_prodotti_in_moncomepos}
	\Todo{}
\end{example}
\begin{example}[Esempi di categorie senza prodotti]\label{ex_cats_senza_prodotti}
	\Todo{}
\end{example}
\Todo{}
%\Todo{Sottolineare che in ciascuna di queste il prodotto soddisfa la stessa proprietà universale che sarà enunciata con precisione in \ref{}}
\subsection{Somme (coprodotti) in \(\ctSet\)}\label{sum_in_set}
\Todo{Qui tutto avviene in maniera duale e in parallelo a prima; cerchiamo di seguire una struttura simile a quella della sezione sui prodotti, per evidenziare la dualità. Però lasciamo molte delle osservazioni pedanti a chi legge, e ci limitiamo a mostrare che la struttura generale della teoria si dà in modo perfettamente duale.}
\begin{itemize}
	\item regole di introduzione ed eliminazione della somma in $\ctSet$;
	\item proprietà universale della somma di insiemi
	\item somma unica up to iso
	\item 3 def equivalenti di somma in $\ctSet$; `somma', `unione disgiunta' e `coprodotto' sono a tutti gli effetti sinonimi in $\ctSet$; in altre categorie si preferisce evitare `unione disgiunta' perché esistono $\ctC$ dove i coprodotti non sono `disgiunti' nel senso di \ref{} e questo potrebbe causare confusione
	\item 3 def equivalenti di somma (se esiste) in $\ctC$, ovvero di coprodotto
	\item osservazioni duali sulla nozione generale di somma
	\item coprodotti in categorie concrete
\end{itemize}
\subsection{Somme in altre categorie `concrete'}
\Todo{Negli spazi topologici, nelle categorie di insiemi ordinati, in alcune piccole categorie di funtori... }
\Todo{In categorie algebriche è un casino: perché? Lo capiremo più avanti.}
\Todo{Si può già fare l'esempio di somme in \(\ctMon\) VS somme in \(\ctCat\); in particolare, \(\susp M +^{\ctCat} \susp N\) è molto diverso da \(M +^{\ctMon} N\).}
\Todo{La proprietà universale della somma è sempre la stessa, \emph{ma le somme in cat di strutture sono molto diverse dalle somme in insiemi!}}
\Todo{Distributività del prodotto sulla somma}
\begin{esercizi}
	\item
	\item
	\item
	\item
	\item
\end{esercizi}
